\section{Discussion}

The evaluation demonstrates the feasibility of an end-to-end,
audit-evidence-producing benchmark for vulnerability-host pair prioritization under
capacity constraints: a 30-seed simulation runs deterministically and reproducibly;
the scheduler and append-only hash-chained audit log operate correctly across all
strategies and seeds; and operational quantities are computed from a frozen,
integrity-verified artifact. The results do \textbf{not} establish real-world
superiority of any strategy, do \textbf{not} show that proposed\_full beats EPSS or
random, do \textbf{not} demonstrate autonomous remediation, do \textbf{not} prove compliance,
and do \textbf{not} constitute production validation. A benchmark that can show *when
added host context does not help* is more valuable than one tuned to confirm a
hypothesis: the finding that uncalibrated context-aware weights do not separate
from EPSS under toy inputs sets a clear falsification condition and motivates
calibrated, larger-window evaluation. Three structural points hold independent of
the (inconclusive) method comparison: per-decision audit records and verifiable
hash chains provide a traceable basis for review; capacity constraints and approval
gates materially shape outcomes (here capacity, not policy, dominated KEV breach);
and the vulnerability-host pair is a useful, traceable unit of decision and
explanation. This work complements prior context-aware/learning-based
prioritization such as Deep VULMAN \textbf{[VERIFY: Deep VULMAN source]}, VulRG [VERIFY:
VulRG source], and VulnScore \textbf{[VERIFY: VulnScore source]} by emphasizing a
reproducible, public-sector-shaped benchmark with per-decision audit evidence and a
freeze/verify protocol, rather than a single tuned result; it does not claim to be
first.
